\begin{abstract}
Motion sickness can make productive laptop use uncomfortable for passengers because visually dominant screen content provides little information about vehicle acceleration and turning. This thesis presents MotionCue, a Windows system that renders peripheral visual motion cues synchronized to inertial measurements from a wirelessly connected smartphone. The system uses native Android and iOS sensor companions, a coordinate-free Bluetooth Low Energy protocol, calibrated vehicle axes, and fail-safe neutralization when sensor frames become stale. The planned evaluation separates system validity from user effectiveness. Study A measures cross-platform sensor agreement, end-to-end latency, packet loss, reconnection, and cue timing. Study B uses a within-participant comparison of no cue, a static visual control, and synchronized cues during a laptop reading task in a passenger vehicle. Primary outcomes are the temporal development of Fast Motion Sickness scores; secondary outcomes include task performance, multidimensional symptoms, workload, and visual distraction. At the current draft stage, the Windows engineering prototype and SensorFrame protocol exist, while smartphone transmission and participant studies remain to be completed. This work is positioned as an exploratory human--computer interaction and sensing study, not as a medical treatment claim.
\end{abstract}
